\subsection{数学速查：遇到符号卡住时先看这里}
\phantomsection\label{app:math-tools}
数学运算生疏时，可以按下面五题定位需要补读的内容，也可以直接查表。
正文会结合例子引入矩阵、导数和概率计算；三角函数、指数与对数也可在用到时回查。

\Needspace{8\baselineskip}
\subsubsection*{五题数学自检：按需定位补读内容}
\phantomsection\label{app:math-entry}
先独立试做，再核对下一段答案。不会某一题，也可以先读正文，遇到相应运算再回来补读。
\begin{enumerate}[itemsep=1pt,topsep=2pt]
\item 解 $2w+1=5$。
\item 计算 $(-2)^2+1$，并说明它与 $-2^2+1$ 的区别。
\item 对 $f(x)=2x+1$，求 $f(-1),f(0),f(1)$。
\item 计算 $\ln(e^2)$，并求满足 $e^u=1$ 的 $u$。
\item 对 $g(w)=(w-2)^2$，求 $g'(w)$ 与 $g'(1)$；在 $w=1$ 时，应增大还是减小 $w$ 才先靠近最低点？
\end{enumerate}

\subsubsection*{核对结果与补读建议}
\phantomsection\label{answer:math-entry}
做完再看。\textbf{第一步：}先代入或整理等式，再判断运算顺序。\par\nobreak
\textbf{结果：}1. $w=(5-1)/2=2$；2. $5$ 与 $-3$，先算幂再应用前面的负号；
3. $(-1,1,3)$；4. $2$ 与 $u=0$；
5. $g'(w)=2(w-2)$，$g'(1)=-2$，导数为负，应先增大 $w$。\par
\textbf{补读：}前 3 题先按答案核对移项、代入与运算顺序；第 4 题查下表中的指数与自然对数；第 5 题查本节后面的“导数、链式法则与一次更新”。

\begin{center}
\small
\renewcommand{\arraystretch}{1.36}
\tagpdfsetup{table/header-rows={1}}
\begin{tabularx}{\linewidth}{@{}p{31mm}X@{}}
\toprule
符号 & 最简单的意思与算例\\
\midrule
$\sum_{i=1}^3 a_i$ & 三项相加：$a_1+a_2+a_3$\\
$\exp(t)=e^t$ & 以自然常数 $e\approx2.71828$ 为底的指数；$e^0=1$\\
$\ln x$ & 自然对数，以 $e$ 为底；$\ln(e)=1$，$\ln1=0$\\
$f'(x),\ \frac{\mathrm df}{\mathrm dx}$ & 小幅改变 $x$ 时，函数输出变化的局部比率\\
$\frac{\partial L}{\partial w}$ & 多个变量中只改变 $w$ 时的局部变化率\\
$\nabla_{\!\theta}L$ & 各参数偏导构成的向量或张量集合\\
$\cos\theta$ & 两条向量夹角的余弦；余弦定理属高中三角知识\\
$\R^{T\times d}$ & $T$ 行 $d$ 列矩阵，本书主线每一行代表一个 Token\\
$\top$ & 转置：行与列交换，$(AB)^\top=B^\top A^\top$\\
$\odot$ & 相同位置的元素逐一相乘，不是矩阵乘法\\
$\approx$ & 数值近似而非精确相等，例如 $1/3\approx0.333$\\
\bottomrule
\end{tabularx}
\end{center}

\subsubsection*{导数、链式法则与一次更新}

如果 $f(x)=x^2$，导数为 $f'(x)=2x$；
例如 $x=3$ 时斜率为 $6$。
如果 $L=(2w-5)^2/2$，令 $u=2w-5$，则
\[
\frac{\mathrm dL}{\mathrm dw}
=\frac{\mathrm dL}{\mathrm du}\frac{\mathrm du}{\mathrm dw}
=u\cdot2=2(2w-5).
\]
这就是\textbf{链式法则}：沿计算路径把局部导数相乘。
学习率 $\eta$ 的梯度下降一步是
$w_{\rm new}=w-\eta\,\mathrm dL/\mathrm dw$。

\subsubsection*{指数与对数怎样求导？}
\label{app:exp-log-derivatives}

Softmax 与交叉熵需要两条求导规则：
\[
\frac{\mathrm d e^z}{\mathrm dz}=e^z,\qquad
\frac{\mathrm d\ln u}{\mathrm du}=\frac1u\quad(u>0).
\]
若对数里面还有计算，要再乘内层的导数。
例如 $f(z)=\ln(1+e^z)$，外层给出 $1/(1+e^z)$，内层给出 $e^z$，所以
$f'(z)=e^z/(1+e^z)$；在 $z=0$ 时为 $1/2$。
分类损失章将逐项用这两条规则求出 $p-y$。

\subsubsection*{均值、方差与独立波动相加}
\phantomsection\label{app:variance-scaling}

均值描述中心，标准差描述数值在中心附近的典型波动幅度。
先看两个独立的数，各以一半概率取 $+1$ 或 $-1$。
它们相加的四种等概率结果是 $(-2,0,0,2)$，均值为 $0$。
\textbf{方差}先把每个结果减去均值，再平方、取平均；
\textbf{标准差}是方差的平方根：
\[
\begin{aligned}
\mu&=\frac{-2+0+0+2}{4}=0,\\
\sigma^2&=\frac{(-2-0)^2+(0-0)^2+(0-0)^2+(2-0)^2}{4}=2,
\qquad \sigma=\sqrt2.
\end{aligned}
\]
这里计算的是四种等概率结果的总体方差，分母为 $4$。
每个原始数的方差是 $1$；两个数独立，和的方差就是 $1+1=2$。
一般地，独立项相加时方差相加，标准差再开平方。

缩放点积的简化假设是：Q/K 的所有坐标相互独立，且各以一半概率取 $\pm1$。
于是每项 $q_i k_i$ 也等概率取 $\pm1$，均值 $0$、方差 $1$；
$d_k$ 项相加得到点积，方差为 $d_k$、标准差为 $\sqrt{d_k}$。
将点积除以 $\sqrt{d_k}$，方差就变为 $1$。
这控制的是波动尺度，数值仍可能超出 $[-1,1]$。

\subsubsection*{一张对照表解决 Token ID 和矩阵下标}

教学词表 ID、位置与数学行列都从 $1$ 起，直接按编号查表：
\[
\begin{array}{c|ccc}
\text{词表 ID／数学列号} & 1 & 4 & 5\\\hline
\text{Python 列下标} & 0 & 3 & 4
\end{array}
\]
“花”（ID 4）的第一位置概率是 $\Pi_{1,4}$。
写 Python 时，在数组索引处统一减 $1$，读取 \texttt{probs[0,3]}。
脚本保留从 $1$ 起的 ID，并将 \texttt{ids - 1} 用于查表、
\texttt{targets - 1} 用于交叉熵标签。

\subsubsection*{不要把矩阵的行列和向量方向混用}

例如 $\vect x^\top=(1,2)$ 作为\textbf{行向量}，
$W=\begin{bmatrix}1&0\\0&3\end{bmatrix}$，
则 $\vect x^\top W=(1,6)$；如果把
$\vect x=(1,2)^\top$ 改写成列向量，同一映射写为
$W^\top\vect x=(1,6)^\top$。
权重矩阵也要转置，不能只把输入横过来。


\par\smallskip
先认清：每个符号是什么、尺寸能否相乘、数字能否代入核对。
到梯度和 Softmax 时，你无需预先记住一整套高等数学；
每遇到一个新运算，都回到这里的小例子。
